\chapter*{Abstract}
\addcontentsline{toc}{chapter}{Abstract}

Object segmentation is an important technique in visual data analysis that identifies and localizes objects at the mask level, and is therefore useful for traffic monitoring, urban surveillance, emergency response, and intelligent transportation systems. This work highlights the application of YOLO-based segmentation in EO/IR aerial traffic imagery, where traffic objects are captured from elevated viewpoints under visible and thermal sensing conditions. Electro-optical (EO) imagery provides rich color and texture information, while infrared (IR) imagery provides thermal cues useful in low-light or visually degraded conditions. However, both modalities differ significantly in appearance, contrast, texture, and boundary structure, creating a domain gap for deep learning based segmentation models. Several training strategies such as EO-only learning, grayscale-domain transformation, supervised IR learning, mixed EO+IR learning, model scaling, high-resolution training, and mask-aware ensembling were explored to improve IR-domain segmentation. The proposed methodology adopts a controlled experimental framework in which each stage studies one major factor affecting segmentation performance. The results show that direct EO-to-IR transfer gives weak performance and grayscale transformations provide limited improvement. The largest gain is obtained when real IR supervision is included during training. Larger models and high-resolution training further improve strict mask quality, while the evaluated ensemble does not outperform the best high-resolution single model. Overall, the study shows that robust EO/IR aerial segmentation depends strongly on target-domain supervision, mixed-modality training, model capacity, and preservation of spatial detail during training.

\vspace{10pt}
\noindent\textbf{Keyword :} EO/IR aerial imagery, infrared segmentation, YOLO segmentation, domain adaptation, mixed-modality training, high-resolution training.
